
\paragraph{\textbf{Remark 2.12.}}
\SGARefTarget{sga3:xiv:2.12}
I do not know whether the assertions of \SGARefLink{sga3:xiv:2.11}{2.11}~a) and b) remain valid
without the hypothesis that regular points exist, when \(k(s)\) is
finite.

It remains to prove
\[
  (\mathrm C_3)\Longrightarrow(\mathrm C_2).
\]
We also record the following equivalent form of
\((\mathrm C_3)\):

\smallskip
\noindent
\textup{\((\mathrm C'_3)\)}
Condition \((\mathrm C_0)\) holds, that is, the nilpotent rank of
\(\mathfrak g(s)\), for \(s\in S\), is locally constant; and on
every open subset \(V\) of \(S\) on which this rank has the value
\(r\), the Killing polynomial of \(\mathfrak g|_V\) is divisible
by \(t^r\).

\smallskip
We must show that this condition implies that every quasi-regular
section \(a\) of \(\mathfrak g\) is regular. In view of \SGARefLink{sga3:xiv:2.6}{2.6}, this
follows from the implication
\(\textup{(iv)}\Longrightarrow\textup{(iii)}\) in the following
lemma, applied to the endomorphism
\(\operatorname{ad}(a)\) of \(\mathfrak g\).

\paragraph{\textbf{Lemma 2.13.}}
\SGARefTarget{sga3:xiv:2.13}
{\itshape Let \(A\) be a ring, let \(M\) be a projective
\(A\)-module of finite type, and let \(u\) be an endomorphism of
\(M\). The following conditions are equivalent:

\smallskip
\noindent
\textup{(i)}
\(M\) is the direct sum of two \(u\)-stable submodules \(M'\) and
\(M''\), such that \(u|_{M'}\) is nilpotent and \(u|_{M''}\) is an
automorphism of \(M''\).

\smallskip
\noindent
\textup{(ii)}
There exists an integer \(n>0\) such that
\[
  \operatorname{Im}u^n+\operatorname{Ker}u^n=M.
\]

\smallskip
\noindent
\textup{(iii)}
The nilspace
\[
  \Nil(u)=\bigcup_{n>0}\operatorname{Ker}u^n
\]
is a direct factor of \(M\), and
\[
  M=\Nil(u)+u(M).
\]

\smallskip
These conditions are implied by the following one, and are
equivalent to it when \(A\) is reduced:

\smallskip
\noindent
\textup{(iv)}
Locally on \(\operatorname{Spec}(A)\), for the Zariski topology, the
characteristic polynomial \(P_u(t)\) of \(u\) can be written in the
form
\[
  t^r\bigl(t^{n-r}+c_1t^{n-r-1}+\cdots+c_{n-r}\bigr),
\]
where \(c_{n-r}\) is invertible.}

The equivalence of \textup{(i)}, \textup{(ii)}, and \textup{(iii)}
is immediate and is recalled for reference. When \(A\) is reduced,
\textup{(i)} implies \textup{(iv)} because a nilpotent
endomorphism of a projective module of rank \(r\) then has
characteristic polynomial \(t^r\). On the other hand, in all cases
the constant term of the characteristic polynomial of an
automorphism of a projective module of finite type is, up to sign,
the determinant of \(u\), locally on \(\operatorname{Spec}(A)\);
it is therefore an invertible element of \(A\).

Finally, to prove
\(\textup{(iv)}\Longrightarrow\textup{(i)}\), regard \(M\) as a
module over the polynomial ring \(A[t]\), with \(t\) acting as
\(u\). The familiar identity
\[
  P(u)=0
\]
shows that \(M\) is annihilated by \(P A[t]\), and hence may be
regarded as a module over \(A[t]/P A[t]\). Write
\[
  P=t^rQ,
\]
where the constant term of \(Q\) is invertible. Then
\[
  P A[t]=t^rA[t]\cap Q A[t],
\]
so that \(A[t]/P A[t]\) decomposes as the product of the rings
\[
  A[t]/t^rA[t]
  \qquad\text{and}\qquad
  A[t]/Q A[t].
\]
This gives a corresponding decomposition of \(M\) into the direct
sum of two \(A[t]\)-modules, or equivalently two \(u\)-stable
\(A\)-submodules \(M'\) and \(M''\); this is precisely the
decomposition of \textup{(i)}.

This completes the proof of \SGARefLink{sga3:xiv:2.13}{2.13}, and hence that of \SGARefLink{sga3:xiv:2.9}{2.9}.
